\originalchapter{曲面与欧拉示性数}\label{ch:surfaces}
\emph{有限复形、边界矩阵、可定向性及角亏依据Paul Seidel的18.900第29--31、33、40讲; 胞腔模型参照18.905第14、18讲. 第40讲为公开补充讲义, Spring 2023课堂未讲. 下文补齐略证并纠正符号笔误; 多边形模型及练习为编者补充.}

\originalsection{有限复形与曲面的局部条件}
多边形粘合能制造曲面, 但任意粘合不一定仍局部像平面. 三角剖分把这个要求转化为顶点周围的三角形如何排列.
\begin{definition}
有限二维单纯复形$K$由有限顶点集、无向边和三角形组成, 每条边的端点、每个三角形的边和顶点都属于$K$, 同一顶点集只允许一个单形. 几何实现$|K|$可取标准基向量为顶点, 将每个单形实现为它们的凸包; 两个单形仅在共同面相交. 若每条边恰邻接两个三角形, 且每个顶点的链环为一个圈, 称为闭组合曲面.
\end{definition}
顶点$v$的链环由与$v$共同成边的邻点, 以及与$v$共同成三角形的邻点对组成. 要求链环为圈意味着围绕$v$的三角形形成一圈扇面, 而不是两圈扇面仅共用$v$. 在边内部, 两个半三角形拼成圆盘邻域; 在顶点处, 圈的锥也是圆盘; 面内部本来就是平面开集. 因而$|K|$是紧拓扑曲面. 反例是把两个三角剖分球面仅在一点粘合: 每条边仍邻接两面, 但粘合点的链环为两个不交圈, 不是流形点.

有边界的组合曲面允许边邻接一个或两个三角形, 顶点链环为区间或圈. 区间链环处得到半圆盘邻域, 构成曲面的边界. 本章讨论有限组合模型, 不声称证明了任意拓扑曲面的可三角剖分性.

\begin{definition}
有限二维复形有$V$个顶点、$E$条边、$F$个三角形时, 定义$\chi(K)=V-E+F$, 称为欧拉示性数. 更一般的有限胞腔分解按各维胞腔数取交错和; 对曲面的多边形分解仍为$V-E+F$.
\end{definition}
细分一条边时增加一个顶点和一条边, 若它属于三角形, 每个邻面又同时增加一条边与一个三角形, 所以$\chi$不变. 在三角形内部加一个顶点并连到三个角时, 变化为$\Delta V=1$, $\Delta E=3$, $\Delta F=2$, 也不改变$\chi$. 更一般, 在一个有$m$条边的面内加一中心并扇形细分, 变化$1-m+(m-1)=0$. 因而这里使用的边细分与面细分均保持$\chi$; 存在共同此类细分的两分解给出同一数.

欧拉示性数对所有有限CW空间的同伦不变性是18.905第18讲的更强结论, 其证明调用前面建立的同调理论. 本册下面证明有限链复形的Euler--Poincar\'e恒等式, 不把它单独冒称为完整的奇异同调不变性证明, 也不据细分计算宣布已证明所有曲面的分类.

\originalsection{边界矩阵与Betti数}
给顶点编号, 每条边任选方向, 每个三角形任选循环方向. 以顶点、边、三角形为基, 定义实向量空间$C_0,C_1,C_2$. 对有向边$[i,j]$令$\partial_1[i,j]=[j]-[i]$; 对有向面$[i,j,k]$令
$\partial_2[i,j,k]=[j,k]-[i,k]+[i,j]$,
反向边按负号处理. 矩阵分别记$D_1,D_2$. 边的方向与顺序在两矩阵中必须一致, 不然乘积没有所声称的意义.

\begin{proposition}\label{thm:boundary-square}
$\partial_1\partial_2=0$. 改变边、面的方向或排序不改变下述同调空间的维数.
\end{proposition}
\begin{proof}
对单个面展开, 得$([k]-[j])-([k]-[i])+([j]-[i])=0$, 线性性给出一般结论. 方向改变是对应基向量乘$-1$, 排序是基置换, 均为可逆基变换; 同一个边基变换同时改变$D_1$的列与$D_2$的行, 核、像和商空间维数不变.
\end{proof}

\begin{definition}
定义$H_0=C_0/\operatorname{im}\partial_1$, $H_1=\ker\partial_1/\operatorname{im}\partial_2$, $H_2=\ker\partial_2$. $b_i=\dim_\R H_i$称为实系数Betti数. 商$H_1$良定义, 正因$\operatorname{im}\partial_2\subseteq\ker\partial_1$.
\end{definition}
这里“闭链”指边界为零的链, “边界链”指某个更高维链的边界. 闭环的边向量是闭链, 因为各中间端点消去. 三角形的边界是可以在面内填起的闭链. $H_1$记录未被面边界消去的线性环绕信息, 会失去基本群的非交换次序.

\begin{theorem}[Euler--Poincar\'e恒等式]\label{thm:euler-poincare}
有限二维复形满足
$b_0=V-\operatorname{rank}D_1$,
$b_1=E-\operatorname{rank}D_1-\operatorname{rank}D_2$,
$b_2=F-\operatorname{rank}D_2$,
故$\chi=b_0-b_1+b_2$.
对任意有限链复形, 链空间维数的交错和等于同调维数的交错和.
\end{theorem}
\begin{proof}
秩零度定理给出$\dim\ker\partial_1=E-\operatorname{rank}D_1$, 再减去其子空间$\operatorname{im}\partial_2$的维数得到$b_1$. 另两式直接由商维数与秩零度定理得到, 相减后秩项消去. 一般有限链复形令$r_i=\operatorname{rank}\partial_i$, 首尾缺失边界的秩取零, 则$b_i=\dim C_i-r_i-r_{i+1}$. 乘$(-1)^i$求和时每个$r_i$恰消去, 得一般恒等式.
\end{proof}

\begin{proposition}
$b_0$等于有限复形的连通分支数.
\end{proposition}
\begin{proof}
每条边的边界将它的两个端点在$H_0$中视为相等. 同一图分支内各顶点沿边路相等, 不同分支之间没有边界关系. 给每个分支取一个顶点类, 它们生成$H_0$且独立: 对每个分支求该分支的顶点系数和, 所得线性泛函在每条边的边界上为零, 并分辨这些类. 面不会连接新的顶点分支, 每个图分支的几何实现道路连通且相互不交, 有限性使它们同时开闭, 所以就是空间的连通分支.
\end{proof}

\begin{example}
计算四面体边界的$\chi$与三个Betti数.
\end{example}
\begin{solution}
$V=4$, $E=6$, $F=4$, $\chi=2$. 边界$D_1$的像由$[2]-[1],[3]-[1],[4]-[1]$生成且独立, 秩为$3$. 将各面方向选成外向, 四个面边界的和为零. 若一组面系数的边界为零, 每条边两侧面系数相等, 四个面相互邻接, 因而全部相等. 所以$\ker D_2$一维, $\operatorname{rank}D_2=3$. 得$(b_0,b_1,b_2)=(1,0,1)$.
\end{solution}

\originalsection{可定向性}
\begin{definition}
若可给各三角形选择方向, 使共边上两面的诱导方向相反, 称闭组合曲面可定向. 否则称不可定向. 对带边界曲面, 条件只在邻接两面的边上检查.
\end{definition}
定向不是把所有平面画出的三角形一律画成逆时针: 图中被分开画出的配对边仍须按实际粘合检查. 一个圈走一周后局部方向可能翻转, 正是M\"obius带的现象.

\begin{lemma}\label{lem:dual-connected}
连通闭组合曲面的面邻接图连通, 其中两面共边时连边.
\end{lemma}
\begin{proof}
以共边可达性划分面. 同一顶点周围的面沿链环依次共边, 故都在同一类. 因而不同类既不共边也不共顶点. 每类的面连同其边、顶点形成子复形, 各类互不相交; 有限个闭子复形的并覆盖曲面, 每类补集也闭, 所以每类同时开闭. 连通性只允许一类.
\end{proof}

\begin{theorem}\label{thm:surface-b2}
连通闭组合曲面在实数系数下, 可定向时$b_2=1$, 不可定向时$b_2=0$.
\end{theorem}
\begin{proof}
选任意初始面方向. 闭$2$链$c=\sum_T a_T T$沿每条边的系数为零. 两侧面系数因此满足$a_{T'}=\pm a_T$, 正负号由面方向决定. 面邻接图连通, 所以从一个面的系数决定全部, 核维数至多$1$. 若曲面可定向, 按相容方向选基, 所有$a_T=1$给出非零闭链, 故核维数恰$1$.

反向若有非零闭链, 传播关系使每个$a_T$非零且绝对值相同. 把负系数面的方向反转, 使所有系数相同且正, 此时每条边的零边界条件迫使两侧诱导方向相反, 给出定向. 因而不可定向时核为零.
\end{proof}
这里底域条件有实质作用. 在$\mathbb F_2$中正负号相同, 所有面的和总是闭链, 连通闭曲面的$b_2$均为$1$. 因而$\mathbb RP^2$的实系数$b_2=0$与模$2$系数$b_2=1$不矛盾.

\begin{proposition}\label{thm:orientable-even}
闭可定向组合曲面的$\chi$为偶数. 连通时$b_1$也是偶数.
\end{proposition}
\begin{proof}
用相容方向把每个面的三条边看作有向“面侧”, 集合大小$3F=2E$. 置换$\tau$将每个面侧送到相邻面的同一边, 因方向相反而翻转, 是$E$个换位, 符号为$(-1)^E$. 置换$\sigma$在同一面中将面侧送到下一条, 为$F$个三循环, 符号为$1$.

$\sigma\tau$保持有向边的起点: 先过到对面的反向边, 再取该面下一边, 起点回到原起点. 顶点链环条件使它在每个顶点周围形成一个循环, 恰有$V$个循环. 一个有$N$个元素、$c$个循环的置换符号为$(-1)^{N-c}$, 故$\operatorname{sign}(\sigma\tau)=(-1)^{2E-V}=(-1)^V$. 同时它等于$(-1)^E$, 所以$V-E$为偶数. $3F=2E$还使$F$为偶数, 因而$\chi=V-E+F$为偶数. 连通可定向时$b_0=b_2=1$, Euler--Poincar\'e式给$b_1=2-\chi$, 亦为偶数.
\end{proof}

\originalsection{多边形曲面模型}
将一个多边形的同名边按箭头方向配对, 可得到一个有限胞腔曲面. 沿多边形边界逆时针读边, 同向记字母, 反向记逆字母. 这是一种胞腔分解, 允许同边两端识别, 不直接满足单纯复形“同顶点集只一个单形”的限制. 可以具体细分为单纯三角剖分. 先将每对边同步三等分, 使每个商边上有两个不同内点, 即使原边两端识别也不产生环边或重边. 在尚未粘合的多边形内, 取一圈不作任何边界识别的互异顶点$u_i$, 分别对应逐段边界顶点$b_i$. 将环带四边形$b_i b_{i+1}u_{i+1}u_i$沿$b_i u_{i+1}$分成两面, 再以一个内部中心$o$扇形细分内圈. 每个三角形有不同的三个顶点, 含有不被识别的内圈顶点, 不同面也不具有同一顶点集, 所以得到真正的单纯复形. 原边三等分不改变$V-E$. 若细分后边界有$p$段, 后续插入$p+1$个顶点及$4p$条内边, 用$3p$面替换原一面, 变化$\Delta\chi=(p+1)-4p+(3p-1)=0$. 这给出胞腔计数与三角剖分计数的一致性, 两种记账仍须区分.

\begin{center}
\begin{tikzpicture}[x=1cm,y=1cm,>=stealth]
\draw (0,0) rectangle (2,1.5);
\draw[->] (.55,0)--(1.45,0); \node[below] at (1,0) {$a$};
\draw[->] (.55,1.5)--(1.45,1.5); \node[above] at (1,1.5) {$a$};
\draw[->] (0,.35)--(0,1.15); \node[left] at (0,.75) {$b$};
\draw[->] (2,.35)--(2,1.15); \node[right] at (2,.75) {$b$};
\node at (1,-.8) {环面};
\draw (4,0) rectangle (6,1.5);
\draw[->] (4.55,0)--(5.45,0); \node[below] at (5,0) {$a$};
\draw[->] (4.55,1.5)--(5.45,1.5); \node[above] at (5,1.5) {$a$};
\draw[->] (4,.35)--(4,1.15); \node[left] at (4,.75) {$b$};
\draw[->] (6,1.15)--(6,.35); \node[right] at (6,.75) {$b$};
\node at (5,-.8) {Klein瓶};
\end{tikzpicture}
\end{center}
\noindent 图中只配对同名边, 箭头同向端点对应. 左图边界词为$aba^{-1}b^{-1}$, 右图为$ab^{-1}a^{-1}b^{-1}$; 将右图的$b$换名为$b^{-1}$, 即得到上一章的$aba^{-1}b$约定.

\begin{example}
计算球面、环面、实射影平面、Klein瓶的胞腔数、基本群及实系数Betti数.
\end{example}
\begin{solution}
球面可将圆盘的整条边界缩成一点, 得一个$0$胞腔和一个$2$胞腔, $\chi=2$. 基本群平凡, 三角模型的Betti数为$(1,0,1)$. 环面有一个顶点、两条边、一个面, $\chi=0$, 基本群$\mathbb Z^2$, 可定向, 因而$(b_0,b_1,b_2)=(1,2,1)$.

实射影平面用圆盘边界对径粘合, 有一个顶点、一条边、一个面, $\chi=1$, 基本群$\mathbb Z/2\mathbb Z$. 它不可定向: 沿射影直线一周, 对径识别使横向方向反转; 也可在圆盘的两条配对半边上检查, 同一面会要求相同方向的边互相消去, 无法选定相容方向. 实系数$b_2=0$, 得$(1,0,0)$. Klein瓶有一个顶点、两条边、一个面, $\chi=0$, 群为上一章的非交换群. 反向粘合的一对边同样翻转横向方向, 不可定向, 得$(1,1,0)$.

上述胞腔计数可由细分多边形验证: 在边上和面中添加的顶点、边、面按前述细分公式成对消去. 球面的两半球三角剖分、环面和Klein瓶的方格细分、射影圆盘边界细分均给同样数值. 有限扭元素不能由实系数Betti数看见, 所以$\mathbb RP^2$的$b_1=0$并不意味着基本群平凡.
\end{solution}

\begin{definition}
连通和$M\#N$通过在两闭曲面各删去一个小开圆盘, 再将新边界圆周配对得到. 定向曲面采用使边界方向相反的配对. 记$\Sigma_g$为$g$个环面的连通和, $g=0$时取球面; 记$N_k$为$k$个实射影平面的连通和, $k\ge1$.
\end{definition}
\begin{proposition}\label{thm:connected-sum-euler}
对上述有限组合模型,$\chi(M\#N)=\chi(M)+\chi(N)-2$. 因而$\chi(\Sigma_g)=2-2g$, $\chi(N_k)=2-k$. 相应基本群表示为
\[
 \pi_1(\Sigma_g)=\left\langle a_1,b_1,\ldots,a_g,b_g\ \middle|\
 \prod_{i=1}^g a_ib_ia_i^{-1}b_i^{-1}=1\right\rangle,
 \quad
 \pi_1(N_k)=\langle c_1,\ldots,c_k\mid c_1^2\cdots c_k^2=1\rangle.
\]
第一式$g=0$时群平凡.
\end{proposition}
\begin{proof}
先细分使删去的盘为一个面. 每删去一个开面使$F$减$1$, 边界仍保留. 两边界均细分为$m$条边后粘合, $V$与$E$各减少$m$, 所以粘合不改变两穿孔曲面的欧拉数之和, 得$-2$公式. 环面数值为零, 射影平面数值为一, 归纳得到两个公式.

将各穿孔多边形从孔边界连一条割线到原边界, 沿割线切开, 得原边界词与孔边界词组合的盘. 连通和粘合孔边界时, 两条孔词方向相反而消去, 保留两个原边界词的顺序拼接. 反复进行后, $\Sigma_g$用一个$4g$边形、边界词为各交换子的积; $N_k$用一个$2k$边形、边界词为各平方的积. 所有多边形角点在这些配对下成为一个顶点: 对每个交换子块四角依次识别, 对每个平方块两边首尾识别, 相邻块的角点连接同一类. 一维骨架于是为相应数目的圆周楔和, 用附加胞腔定理得到群表示.
\end{proof}
这些模型给出重要曲面族及其不变量, 不等于已经证明“每个紧连通曲面都在此表中”. 曲面分类还须证明一般多边形词可化为标准形, 并核对唯一性; 任意拓扑曲面还涉及三角剖分定理, 均不在本册证明范围.

\originalsection{角亏与离散Gauss--Bonnet}
给每条边指定正长度, 每个三角形的三边满足严格三角不等式, 就能按该长度实现为欧氏三角形并沿等长边粘合. 这里的内在几何不要求整个曲面嵌入$\R^3$, 相邻面可以在空间中没有指定的夹角.
\begin{definition}
闭组合曲面的顶点$v$处, 定义角亏$\kappa_v=2\pi-\sum_{T\ni v}\theta_{T,v}$, 其中$\theta_{T,v}$为面$T$在$v$处的内角. 正角亏表示总角小于一整周, 负角亏表示总角大于一整周.
\end{definition}
\begin{theorem}[离散Gauss--Bonnet]\label{thm:discrete-gauss-bonnet}
有限闭欧氏三角曲面满足$\sum_v\kappa_v=2\pi\chi$. 有边界时, 边界顶点改用$\kappa_v=\pi-\sum_{T\ni v}\theta_{T,v}$, 仍有$\sum_v\kappa_v=2\pi\chi$.
\end{theorem}
\begin{proof}
闭情形每面三角和为$\pi$, 故$\sum_v\kappa_v=2\pi V-\pi F$. 每条边邻接两面, 数面边关联得到$3F=2E$, 因而$2\pi(V-E+F)=2\pi V-\pi F$.

有边界时记边界顶点数与边数为$V_b,E_b$, 内部顶点数为$V_i$. 边界是有限个组合圆周, 故$V_b=E_b$; 关联计数给$3F=2E-E_b$. 于是
\[
 \sum_v\kappa_v=2\pi V_i+\pi V_b-\pi F
 =2\pi(V_i+V_b-E+F)=2\pi\chi.
\]
无需曲面可定向; 两种计数只依赖局部曲面条件.
\end{proof}

\begin{example}
求正二十面体各顶点的角亏, 并说明平坦环面为何可以处处零角亏.
\end{example}
\begin{solution}
每顶点邻接五个等边三角形, $\kappa=2\pi-5\pi/3=\pi/3$, 十二顶点总和$4\pi=2\pi\cdot2$. 平面等边三角形格按两个独立平移周期取商, 每点邻接六个$\pi/3$角, 角亏为零, 符合环面$\chi=0$. 若需严格单纯复形, 取足够大的周期格避免重复边和相同顶点面的识别; 小周期图可仅作为胞腔记账, 不能冒称单纯剖分.
\end{solution}
角亏总和为零不要求每一点都零: 可以有正负角亏互相抵消. 离散定理不等于光滑曲面的高斯曲率积分定理, 后者需要新的曲率定义与分析论证.

\originalsection{练习与解答}
\begin{exercise}
在有四个顶点的复形中, 取所有六条边, 但只填三个三角面$[1,2,3]$, $[1,2,4]$, $[1,3,4]$. 求Betti数, 并说明是否为闭曲面.
\end{exercise}
\begin{solution}
$\operatorname{rank}D_1=3$. 三面的边界独立, 因为边$[2,3],[2,4],[3,4]$各只属于其中一个面, 零边界迫使三个系数都零, 故$\operatorname{rank}D_2=3$. 得$(b_0,b_1,b_2)=(1,0,0)$, $\chi=1$. 它为四面体边界删去一个开面的圆盘, 三条外边只邻接一个面, 故不是闭曲面.
\end{solution}
\begin{exercise}
设闭三角曲面全部面为等边三角形, 每个顶点恰有六条邻边. 求$\chi$. 能否据此断言它为环面?
\end{exercise}
\begin{solution}
链环为圈使邻面数等于邻边数, 每顶点总角为$2\pi$, 角亏为零, 故$\chi=0$. 仅此不能断言环面, 因为Klein瓶也有这种平坦组合模型, 且不要求曲面连通. 即使另外假设连通可定向, 从$\chi=0$推出环面还要调用本册未证明的完整曲面分类定理.
\end{solution}
\begin{exercise}
对边闭链$c$及满足$D_2^{\mathsf T}u=0$的边向量$u$, 证明$u\cdot c$只依赖$c$在$H_1$中的类. 若$u=D_1^{\mathsf T}v$, 证明它对所有闭链均为零.
\end{exercise}
\begin{solution}
换代表$c'=c+D_2w$时, $u\cdot(c'-c)=(D_2^{\mathsf T}u)\cdot w=0$. 若$u=D_1^{\mathsf T}v$, 则$u\cdot c=v\cdot D_1c=0$. 这些线性绕数只看边的净次数, 不看排列次序, 所以不能代替基本群; 例如两圈楔和的交换子对所有此类实值测量都为零, 但在自由群中非平凡.
\end{solution}
